\section{Discussion}\label{sec:discussion}

\paragraph{What we learned.} Training a network on an upper theory gave it
the lower theory only in a jagged form, and that jagged form is where its
jagged competence came from. The network did learn something of A: the part
of A that B depends on became readable, while the parts B ignores became less
readable. But
what it learned was uneven in the places that matter for B. Category misreads
concentrated near the slot-1 cutoffs, where B's categories change, even
though a probe could read the slot-1 sum on 97--98\% of boards;
when A was trained as an explicit output it was largely but not exactly
right; and part of the prediction's failure sat where this partial A failed.
Seen this way, jagged competence is not noise on top of an otherwise correct
solution. Uneven acquisition, access and preservation of A explain part of
B's jagged competence; some failures remain unexplained. When B instead needed all five sums
(the all-slots task), all five became readable, which suggests that the
jagged form of A follows what B asks of it; that comparison is historical
and crosses tasks.

\paragraph{Knowing an intermediate is not one thing.} Our second lesson is
that ``the network knows A'' hides five separate facts. A can be readable
without established exact computation, computed without reaching the
prediction, delivered in a form that supports less reliable prediction, or
lost under further training, and an exact, delivered A can still leave
failures in B. Matched comparisons isolate the effects of connection,
encoding and prediction gradients: connecting a learned A to B cut misreads
by up to 17-fold, the same exact A gave fewer misreads as one-hot inputs
than as numerical inputs, and a single B update made an exact A inexact unless the
B gradient was blocked. Having the right information is therefore not enough;
how it reaches the decision, in what form, and whether it survives training
all matter. Nor is more supervision of the intermediate automatically
better: networks also trained to output the sums had more misreads (473--815)
than networks trained on the prediction alone (58--328) in r13's
sampled-target condition (\cref{sec:res-noise}).

\paragraph{Jaggedness has locatable sources.} The known layers let us
distinguish several sources of failure: missing access, missing
preservation, a less usable encoding, noisy targets, and history coverage,
assessed in a historical comparison. Some cases remain unexplained: 19 of the 42 diagnosed
r12 histories that exceed the diagnostic bar, and the residual failures with
exact sums. Average error and probe accuracy each looked healthy while these
failures were present, so neither alone establishes complete, reliable
realization of the task. ``Jagged'' names the phenomenon; this account, not the
word, is the explanation.

\paragraph{Relation to other accounts.} Coverage-based accounts explain
jaggedness by gaps in what a system has seen \citep{gans2026}; history
coverage is one such factor here, but the access, preservation and encoding
effects appear in matched comparisons with identical training data, so
coverage alone does not account for them. Other studies report gaps between
what probes read and what outputs do \citep{basu2026,darade2026}; our
connection experiment shows one way such a gap can cause errors and be
reduced.

\paragraph{What we expect to carry over.} We expect the two measurement
cautions to hold wherever the same measurements are used: low average error
does not establish reliability on constructed histories, and recoverability
does not establish predictive use. Whether learning an upper theory yields a
jagged lower theory in larger systems is a test to run, for example by
checking which parts of a known intermediate become readable under
final-output training, whether a learned intermediate output is used by the
main prediction, and whether downstream training preserves it.
